\documentclass[../main.tex]{subfiles}

\begin{document}

\section{Findings Register}
\label{app:findings-register}

This appendix is the authoritative register of formal findings for this
revision and maintains a separate evidence-and-scoping backlog. It applies the
method in Section~\ref{sec:risk-rubric} to public evidence available through
August 27, 2026. A formal finding is a concrete adverse-event scenario at a
stated boundary; it is not a score assigned to a technology, feature, CVE
count, recommendation, or project. Broad mechanisms that do not yet define an
adverse consequence and affected population appear only in
Table~\ref{tab:findings-evidence-backlog}; they are not formal findings. Unless
a record expressly says otherwise, the assessment did not independently
reproduce the reported condition and did not assess the current installed
population.

Table~\ref{tab:findings-register-summary} records inherent-risk axes and tiers
where evaluated, together with every finding's rating status. It does not imply
that every version or deployment is affected today. A public fix, issue
closure, or reporter statement is noted where relevant, but it is not credited
as current residual-risk reduction without sufficient evidence of the
implemented control and its coverage. An \emph{NE} entry is an evidence or
scoping gap, not a Low-risk conclusion.

Linked recommendation identifiers show intended control coverage only.
Response routing is shown separately in
Table~\ref{tab:findings-response-routing}. Implementation dependencies, effort,
delivery sequence, ownership, acceptance evidence, and residual-risk decisions
do not alter the inherent ratings below and are not populated in this findings
register. The proposed work packages, dependencies, accountable roles, effort,
and acceptance gates are maintained separately in
Appendix~\ref{app:phased-roadmap}; the linked identifiers must not be read as
completed planning records.

\begingroup
\scriptsize
\setlength{\tabcolsep}{2pt}
\renewcommand{\arraystretch}{1.15}
\setlength{\LTleft}{-0.7in}
\setlength{\LTright}{-0.7in}
\begin{longtable}{@{}p{0.14\widetablewidth}p{0.33\widetablewidth}p{0.075\widetablewidth}p{0.075\widetablewidth}p{0.10\widetablewidth}p{0.13\widetablewidth}p{0.09\widetablewidth}@{}}
\caption{Formal findings and report-wide rating status}
\label{tab:findings-register-summary}\\
\hline
\textbf{ID} & \textbf{Scenario shorthand} & \textbf{Dim.} & \textbf{I/E} &
\textbf{Tier} & \textbf{Status} & \textbf{Conf.} \\
\hline
\endfirsthead

\caption[]{Formal findings and report-wide rating status, continued}\\
\hline
\textbf{ID} & \textbf{Scenario shorthand} & \textbf{Dim.} & \textbf{I/E} &
\textbf{Tier} & \textbf{Status} & \textbf{Conf.} \\
\hline
\endhead

\hline
\endfoot

\hline
\endlastfoot

\code{CORE-SEC-001} &
Malformed on-disk metadata reaches memory-unsafe core-library paths in
historically affected releases. & Security & I3/E2--E3 & High & Provisional & Moderate \\

\code{CORE-SEC-002} &
File-controlled work, allocation, traversal, or error paths terminate or
exhaust an affected reader. & Security & I2/E2--E3 & Moderate & Provisional & Moderate \\

\code{CORE-TOOL-001} &
A malicious file triggers the reported \code{h5dump} use-after-free in HDF5
1.14.1-2 and earlier. & Security & I2--I3/E2 & Moderate--High & Provisional & Moderate \\

\code{CORE-PRIV-001} &
Sensitive plaintext metadata is disclosed when a file is shared under
payload-only protection. & Privacy & I2--I3/E2 & Moderate--High & Provisional & Moderate \\

\code{CORE-SAFE-001} &
Storage exhaustion during write or close leaves an incomplete or corrupted
file after the operation reports failure, with no defined recovery or
poisoned-file protocol. & Safety & I2--I3/E2 &
Moderate--High & Provisional & Low \\

\code{CORE-SAFE-002} &
Abrupt process or host termination during write leaves a file that is
corrupted or not reopenable. & Safety & I2--I3/E3 & Moderate--High &
Provisional & Low \\

\code{LONG-SAF-001} &
An archived dataset becomes unreadable because a required non-core filter is
no longer available. & Safety & I2/E2 & Moderate & Provisional & Low \\

\code{APP-SEC-001} &
The affected Hugging Face dataset service followed attacker-selected external
raw-storage paths with worker authority and returned local files. & Security & I3/E4 & Critical & Rated & High \\

\code{APP-SEC-002} &
Affected Keras legacy-HDF5 model loading executes embedded code despite
\code{safe_mode=True}. & Security & I3/E2 & High & Rated & High \\

\code{APP-SEC-003} &
Affected Keras model loading follows attacker-controlled HDF5 external dataset
references and discloses local files. & Security & I3/E3 & High & Rated & High \\

\code{APP-SEC-004} &
Affected Keras weight loading allocates from an extreme file-declared shape and
exhausts memory. & Security & I2/E3 & Moderate & Rated & High \\

\code{APP-SEC-005} &
TensorFlow 2.17.0 loads HDF5 plugins from an uncontrolled search location,
enabling local privilege escalation. & Security & I3/E2 & High & Rated & Moderate \\

\code{IND-SAF-001} &
HDF5 2.0/2.1 exact-width validation rejects a sufficient but nonminimal chunk
encoding permitted by the specification. & Safety & I2/E2 & Moderate & Rated & High \\

\code{IND-SAF-002} &
Older permissive readers accept contradictory PureHDF 2.1.2 chunk metadata
that HDF5 2.0 rejects after upgrade. & Safety & I2/E2 & Moderate & Rated & High \\

\end{longtable}
\endgroup

\subsection{Correspondence with the SSP SIG Registries}
\label{sec:findings-ssp-crosswalk}

Table~\ref{tab:findings-ssp-crosswalk} relates the formal findings and
evidence-backlog records maintained in this appendix to the SSP SIG threat,
hazard, and privacy registries. Three cautions govern its use, and all three are
load-bearing.

First, this is an \textbf{editorial correspondence established by this report}.
Neither project has adopted it, and a row asserts only that the two records
describe a related mechanism. Second, \textbf{it is not a rating comparison}:
the registries score Severity $\times$ Likelihood on a $1$--$25$ scale with five
bands, which Section~\ref{sec:rubric-reconciliation} establishes is not
convertible to or from the tiers used here. Third, \textbf{``no direct
counterpart'' runs in both directions and means neither project is wrong}. The
registries model the library and format; this report additionally rates
application, service, and independent-implementation boundaries that the
registries do not cover, and the registries enumerate operational safety
hazards this report did not examine.

\begingroup
\small
\setlength{\tabcolsep}{3pt}
\renewcommand{\arraystretch}{1.15}
\setlength{\LTleft}{-0.9in}
\setlength{\LTright}{-0.9in}
\begin{longtable}{@{}>{\raggedright\arraybackslash}p{0.16\widetablewidth}>{\raggedright\arraybackslash}p{0.26\widetablewidth}>{\raggedright\arraybackslash}p{0.54\widetablewidth}@{}}
\caption{Editorial correspondence between this register and the SSP SIG
registries}
\label{tab:findings-ssp-crosswalk}\\
\hline
\textbf{This report} & \textbf{SSP SIG entry} & \textbf{Note} \\
\hline
\endfirsthead

\caption[]{Editorial correspondence between this register and the SSP SIG
registries, continued}\\
\hline
\textbf{This report} & \textbf{SSP SIG entry} & \textbf{Note} \\
\hline
\endhead

\hline
\endfoot

\hline
\endlastfoot

\code{CORE-SEC-001} &
\code{ATK-001}--\code{ATK-008} &
The registry decomposes the memory-unsafe parser class into eight mechanism
entries; this report treats them as one shared adverse event with a common
scope and evidence base. The decomposition is the better basis for the
\code{CORE-S1} invariant tranches. \\

\code{CORE-SEC-002} &
\code{ATK-009}, \code{ATK-010}, \code{ATK-012} &
Resource exhaustion, null and uninitialized dereference, and I/O amplification
from pathological chunk layouts. The development-tree assert-only case tracked
under \code{CORE-SEC-003} corresponds to \code{ATK-010} but is not part of this
finding's rated population. \\

\code{CORE-TOOL-001} &
\code{ATK-007} (mechanism only) &
The registry rates lifetime defects as a library class and does not separate
official-tool exposure. This report rates the tool boundary distinctly. \\

\code{CORE-PRIV-001} &
\code{EXP-002}; also \code{HAZ-004} &
The one pairing where both projects rate the same mechanism and land in
comparable bands. \\

\code{CORE-SAFE-001} &
\code{HAZ-005} &
Same scenario. The registry supplies a hazard analysis and rating; this report
rates the bounded scenario Moderate--High, Provisional, with Low confidence;
\code{RM-2D} measures durable-state and recovery outcomes and may change that
assessment. See Table~\ref{tab:haz-crosswalk}. \\

\code{CORE-SEC-003} &
\code{ATK-010} (partial) &
No registry entry addresses build-mode-dependent enforcement directly. \\

\code{CORE-SAFE-002} &
\code{HAZ-001}; also \code{HAZ-002}, \code{HAZ-006}, \code{HAZ-012},
\code{HAZ-015} &
\code{HAZ-001} is the registry's highest-rated safety entry. This report
registers one finding where the registry has five; the durability cluster
around flush semantics, concurrency, and recovery tooling is not separately
rated here. \\

\code{LONG-SAF-001} &
\code{HAZ-014} &
Same missing-filter portability mechanism, assessed here at an explicit
ten-year-or-more archival horizon. VFD and VOL archival dependencies remain
outside this rated record. \\

\code{APP-SEC-001} &
No direct counterpart &
The registries contain no confused-deputy or authority-boundary entry.
\code{EXP-008} and \code{EXP-009} cover external-link disclosure as a privacy
mechanism, which is related but not the same scenario. \\

\code{APP-SEC-002} &
No direct counterpart &
Application-level deserialization layered on HDF5 is outside the registries'
scope. \\

\code{APP-SEC-003} &
\code{EXP-009} (partial) &
External references revealing or reaching unintended targets; the registry
frames it as disclosure, this report as an authority-boundary failure. \\

\code{APP-SEC-004} &
\code{ATK-009} &
Same mechanism, rated here at the application boundary that performs the
allocation. \\

\code{APP-SEC-005} &
\code{ATK-011} &
The clearest methodological disagreement in the pair of documents; see
Table~\ref{tab:ssp-divergence}. \\

\code{FMT-SAF-001} &
No direct counterpart &
The hazard model has an \code{H6} parser-ambiguity family and H5II has
\textbf{I1} parser/spec divergence, but neither registry records the unresolved
normative-artifact question represented by this evidence-backlog item. \\

\code{IND-SAF-001}, \code{IND-SAF-002} &
No direct counterpart &
The registries have a parser-ambiguity hazard family but no
independent-implementation or reader/writer-conformance entries. This is
material this report contributes. \\

\code{EXT-PRIV-001} &
\code{EXP-004}--\code{EXP-007}; also \code{HAZ-003} &
The registry has already enumerated four bounded extension exposures where this
report has one unrated theme. \code{RM-2E} should adopt them as its
enumeration. \\

\code{SUP-SAFE-001} &
\code{ATK-013}, \code{HAZ-014} &
Distribution compromise and missing-filter portability. \code{ATK-013} is a
security scenario this report never developed; \code{RM-2F} should adopt it
rather than re-derive one. \\

\end{longtable}
\endgroup

\subsection{Evidence and Scoping Backlog}

The observations below identify work needed to determine whether additional
formal findings exist. They have no impact or exposure axes, risk tier,
response urgency, or implied Low-risk status.

\begingroup
\small
\setlength{\tabcolsep}{3pt}
\renewcommand{\arraystretch}{1.15}

\setlength{\LTleft}{-0.9in}
\setlength{\LTright}{-0.9in}

\begin{longtable}{@{}p{0.21\widetablewidth}p{0.29\widetablewidth}p{0.44\widetablewidth}@{}}
\caption{Observations requiring a concrete adverse-event scenario}
\label{tab:findings-evidence-backlog}\\
\hline
\textbf{ID} & \textbf{Observation} & \textbf{Evidence needed for a finding} \\
\hline
\endfirsthead

\caption[]{Observations requiring a concrete adverse-event scenario, continued}\\
\hline
\textbf{ID} & \textbf{Observation} & \textbf{Evidence needed for a finding} \\
\hline
\endhead

\hline
\endfoot

\hline
\endlastfoot

\code{CORE-SEC-003} &
Candidate deserializer invariants derived from in-file bytes are masked inside
\code{assert(...)} expressions and therefore may not be enforced in
\code{-DNDEBUG} release builds. Issue~\#87 identifies seven component families
and estimates 50--100 candidates, but does not provide a verified site catalog
or consequence analysis; a fixed develop-source sample corroborates the build
mechanism and representative sites, not their runtime consequences
~\cite{hdf5-pickles-issue-87,hdf5-assert-build-config-2026,
hdf5-assert-chunk-record-2026,hdf5-assert-extensible-array-2026,
hdf5-assert-fractal-heap-2026}.
One v2 B-tree \code{record\_size} case advances the evidence without closing
the supported-release gap: its source-reported reproduction records
\code{SIGFPE} on an \code{-DNDEBUG}, release-configured HDF5 2.3.0 development
tree, but explicitly records no 32-bit measurement and no affected supported
version or version range~\cite{h5policy-case-v2btree-recordsize}.
A second candidate is reported for the cache-image LRU rank relation without a
measured release-build consequence~\cite{h5policy-case-mdci-lru}. The remaining
population is unquantified, and this record stays open. &
For each supported source line and build profile, record the source site,
file-byte provenance, non-asserting dominating guard if any, release
reachability, first security-relevant sink, paired debug and \code{-DNDEBUG}
fixture result, affected versions, and consequence. Promote each supported
adverse-event scenario into \code{CORE-SEC-001}, \code{CORE-SEC-002}, or a new
formal finding as the evidence requires; a redundant or unreachable assertion
needs a documented negative disposition. Begin from the measured development-
tree case's sibling audit, which names further unaudited assertions in the same
B-tree header initializer and in fractal-heap initialization. \\

\code{FMT-SAF-001} &
The reference library and the published file-format specification disagree on
whether the group-info message is optional, and no ruling had been recorded at
the cutoff. Issue~\#6409 separately demonstrates the value of an independent
implementation, but not indeterminacy: the published rule settled that case
and the reader was corrected~\cite{hdf5-gh-issue-6409,hdf5-gh-pr-6432,
hdf5-gh-issue-6616}. &
Record an authoritative disposition for issue~\#6616, identify any file or
workflow whose behavior depends on the disputed rule, and measure the resulting
access, integrity, or compatibility consequence. Until both an adverse outcome
and affected population are bounded, this is a format-governance observation,
not a rated safety finding. \\

\code{EXT-PRIV-001} &
Extensions can move data and metadata across storage, cache, logging, network,
and executable-plugin boundaries. &
Identify a deployment, sensitive-data population, unauthorized recipient or
retention event, affected scope, and resulting harm. \\

\code{SUP-SAFE-001} &
Downstream builds and workflows can disagree about ABI, linkage, filters,
plugins, VFDs, parallel features, or file locking. &
Identify a concrete producer/consumer profile pair, affected population,
failure mode, recoverability, and recurrence evidence. \\

\end{longtable}
\endgroup

\subsection{Response Routing}

The routing below follows the default rule in
Section~\ref{sec:risk-rubric}. Where current residual risk was not evaluated,
it uses inherent risk as the conservative provisional baseline; a tier range is
routed at its upper bound. Routing is conditional on confirming the stated
affected scope and is not a claim that every present deployment needs the same
action. The target horizons begin when an affected deployment or equivalent
exposure is identified. They are assessment recommendations, not
stakeholder-approved delivery commitments, and do not incorporate effort or
dependencies. In particular, Hugging Face reports that it stopped wrongly
processing HDF5 external references, but this assessment did not validate the
implemented control or assign a current residual tier.
Phases~0--2 of Appendix~\ref{app:phased-roadmap} operationalize these clocks;
later structural work neither replaces nor extends them.

\begingroup
\small
\setlength{\tabcolsep}{3pt}
\renewcommand{\arraystretch}{1.15}

\setlength{\LTleft}{-0.9in}
\setlength{\LTright}{-0.9in}

\begin{longtable}{@{}p{0.22\widetablewidth}p{0.29\widetablewidth}p{0.42\widetablewidth}@{}}
\caption{Provisional response urgency and target horizons}
\label{tab:findings-response-routing}\\
\hline
\textbf{Urgency and horizon} & \textbf{Findings} & \textbf{Decision rule for this register} \\
\hline
\endfirsthead

\caption[]{Provisional response urgency and target horizons, continued}\\
\hline
\textbf{Urgency and horizon} & \textbf{Findings} & \textbf{Decision rule for this register} \\
\hline
\endhead

\hline
\endfoot

\hline
\endlastfoot

Immediate---before the next untrusted processing job and no later than 24 hours
after identification & \code{APP-SEC-001} &
Immediately contain any equivalent service that still automatically follows
file-selected external resources with access to sensitive local files. The
historical incident rating is not a current rating of Hugging Face. \\

Near-term---within 30 days & \code{CORE-SEC-001}, \code{CORE-TOOL-001},
\code{CORE-PRIV-001}, \code{CORE-SAFE-001}, \code{CORE-SAFE-002},
\code{APP-SEC-002}, \code{APP-SEC-003}, \code{APP-SEC-005} &
Identify affected releases and exposed deployments, apply available fixes or
containment, minimize disclosed metadata, and verify the trust boundary. For
the two write-side safety findings the Near-term obligation is \emph{tested
operational containment} --- capacity headroom and monitoring, flush and
checkpoint discipline, a rehearsed recovery procedure, and an independent copy
of data that cannot be regenerated. The library changes in
\code{CORE-SAFE-S1} and \code{CORE-SAFE-S2} are structural and follow the
Planned and Phase~3 horizons. \\

Planned---within 90 days or the next supported release or migration window,
whichever is sooner & \code{CORE-SEC-002},
\code{APP-SEC-004}, \code{IND-SAF-001}, \code{IND-SAF-002},
\code{LONG-SAF-001} &
Plan treatment and acceptance testing; expedite where an affected production
pipeline, archive, sensitive-data population, or service is identified.
For \code{LONG-SAF-001}, identify the bounded archive corpus and its filter
dependencies within the clock; longer-term profile and preservation work does
not extend that first deliverable. \\

\end{longtable}
\endgroup

\subsection{Evidence Actions Without Assigned Urgency}
\label{sec:findings-evidence-actions}

Items \code{CORE-SEC-003}, \code{FMT-SAF-001}, \code{EXT-PRIV-001}, and
\code{SUP-SAFE-001}
require the scenario evidence listed in
Table~\ref{tab:findings-evidence-backlog}. These are evidence actions, not a
fifth response-urgency category.

\code{CORE-SAFE-001} appeared in this category in an earlier revision and no
longer does. The cited API behavior and maintainer discussion establish the
control gap, while the unmeasured frequency and durable-state outcome constrain
exposure to E2 and keep confidence Low.

\paragraph{Horizon handling for
\code{LONG-SAF-001}.}
The I2/E2 Moderate rating maps to \emph{Planned} without a departure. Its
ten-year-or-more horizon remains material, so decision rule 7 still forbids
direct ordering against scenarios that state no horizon. The 90-day deliverable
is a bounded corpus and dependency inventory plus an accountable disposition;
it is not a claim that decadal preservation can be completed in 90 days. If an
archive is already failing, that present-tense instance is separately scoped
and routed without the horizon qualifier.

\subsection{Core Library, Tool, and Format Findings}

\subsubsection*{CORE-SEC-001 --- Memory-unsafe processing of malformed metadata}

\noindent\textbf{Scenario.}
An untrusted originator supplies a crafted HDF5 file to a process using a
historically affected HDF5 C library. Insufficiently validated file-derived
sizes, counts, offsets, message state, layout state, filter parameters, or
datatype state reach allocation, copy, object-construction, or cleanup paths
and cause memory corruption, including credible process compromise for the
variants whose primary records support that consequence.

\noindent\textbf{Scope and evidence.}
The assessed population consists only of deployments matching an affected
version-and-path pair represented in Appendix~\ref{app:cve-history-summary} and
processing untrusted files. The H5FS, H5C, H5O, H5T, H5Z, heap,
dataspace/layout, and vector-copy variants are evidence for the shared adverse
event, not separately scored objects. A detailed primary report for HDF5 1.8.16
traces file-controlled array rank to an out-of-bounds heap write and possible
code execution in the consuming application~\cite{talos-hdf5-2016-0176}; the
cited later fixes demonstrate early sanity checks, bounds and overflow
tracking, and controlled rejection replacing unsafe behavior
\cite{cve_hdf5,commit_804f3ba,commit_3914bb7,commit_ce53bc0}. This record
excludes termination-only and resource-consumption variants assigned to
\code{CORE-SEC-002} and the tool-specific UAF assigned to
\code{CORE-TOOL-001}. It does not assert that every CVE had code-execution
impact or that a current HDF5 2.x release contains every historical condition.

\noindent\textbf{Rating basis.}
I3 reflects credible process compromise or major loss of process integrity in
the evidenced variants. E2 applies to interactive or otherwise constrained
file processing; E3 applies where an affected deployment automatically ingests
untrusted files through a common workflow. Both produce a High tier, while
exact version and deployment prevalence remain incomplete. The rating is
therefore \textbf{I3/E2--E3, High, Provisional, Moderate confidence}. Current
release and installed-population residual risk were not evaluated.

\noindent\textbf{Linked recommendations.}
\code{CORE-S0}--\code{CORE-S7} and \code{CORE-S12}; the primary structural
controls are \code{CORE-S1}--\code{CORE-S3}.

\subsubsection*{CORE-SEC-002 --- Uncontrolled work and resource consumption}

\noindent\textbf{Scenario.}
A crafted, truncated, or inconsistent file reaches affected decode, traversal,
arithmetic, allocation, recursion, or lifecycle paths and causes an assertion,
null dereference, divide-by-zero, leak, stack exhaustion, or uncontrolled
resource consumption, interrupting the consuming process or workflow.

\noindent\textbf{Scope and evidence.}
The assessed population consists only of deployments matching an affected
version-and-path pair in Appendix~\ref{app:cve-history-summary} and processing
the triggering file. This record includes termination and resource-consumption
outcomes without evidenced memory corruption; it excludes the memory-unsafe
variants assigned to \code{CORE-SEC-001} and the tool-specific UAF assigned to
\code{CORE-TOOL-001}. Exact fixed versions across the represented paths and the
present installed population are not established. The file-space repair is one
direct example of converting an assertion reached from corrupted metadata into
a controlled error~\cite{commit_804f3ba,cve_hdf5}.

\noindent\textbf{Related development-tree evidence.}
The \code{CORE-SEC-003} backlog contains a source-reported v2 B-tree case in
which an assert-only zero-record-size guard is absent in an \code{-DNDEBUG}
build and the public API path terminates with \code{SIGFPE}. It was measured on
a locally built HDF5 2.3.0 development source tree, not an established
supported release; the record explicitly says that 32-bit behavior and the
affected version range were not measured
~\cite{h5policy-case-v2btree-recordsize}. It corroborates the mechanism and
provides a concrete fixture for \code{RM-1E}, but it is not included in this
formal finding's rated population. A related zero-capacity variant reported
against an older h5py-linked library is likewise uncharacterized and unrated.

\noindent\textbf{Rating basis.}
I2 reflects a meaningful but ordinarily recoverable process or workflow
interruption; a mission-level outage requires a separate deployment-specific
record. E2 applies to interactive or otherwise constrained file processing;
E3 applies where an affected deployment automatically ingests untrusted files
through a common workflow. Both produce a Moderate tier, while incomplete
version prevalence makes the assessment provisional. The rating is
\textbf{I2/E2--E3, Moderate, Provisional, Moderate confidence}. Current release
and installed-population residual risk were not evaluated.

\noindent\textbf{Linked recommendations.}
\code{CORE-S0}--\code{CORE-S8}, \code{CORE-S12}, \code{APP-S1}, and
\code{APP-S2}.

\noindent\textbf{SSP SIG correspondence.}
\code{ATK-009}. The development-tree assertion case in \code{CORE-SEC-003}
corresponds to \code{ATK-010}, but is outside this finding's rated population.
See Section~\ref{sec:rubric-reconciliation} for why the ratings differ.

\subsubsection*{CORE-TOOL-001 --- h5dump use-after-free}

\noindent\textbf{Scenario.}
An attacker supplies a malicious HDF5 file to a user or automated workflow that
invokes an affected \code{h5dump}; the file triggers the reported
\code{H5T__conv_struct} heap use-after-free and may compromise or terminate the
inspection process.

\noindent\textbf{Scope and evidence.}
The HDF Group advisory identifies HDF5 1.14.1-2 and earlier as affected and
reproduces AddressSanitizer-detected termination when the supplied file is
processed by the utility~\cite{cve-2026-34734,hdf5-ghsa-w7v2-9cmr-pwwj}. It
states that practical code-execution exploitability is unknown. This is a tool
finding, not evidence that every libhdf5 consumer reaches the same path.

\noindent\textbf{Rating basis.}
The demonstrated termination supports I2; credible but unverified process
compromise from the UAF supports I3 as the upper bound. E2 reflects the
affected-version, crafted-file, and explicit inspection preconditions. The
rating is \textbf{I2--I3/E2, Moderate--High, Provisional, Moderate confidence}.
Current installed-population residual risk was not evaluated.

\noindent\textbf{Linked recommendations.}
\code{CORE-S6}, \code{CORE-S9}, \code{CORE-S12}, and \code{APP-S2}.

\subsubsection*{CORE-PRIV-001 --- Plaintext metadata disclosure}

\noindent\textbf{Scenario.}
An HDF5 file containing sensitive object names, attributes, member labels,
comments, user-block content, provenance, or structural relationships is shared
with a recipient who is authorized for a protected payload but not for that
context. Direct HDF5 or byte-level inspection reveals the unintended metadata.

\noindent\textbf{Scope and evidence.}
This finding is limited to files that actually contain sensitive plaintext
metadata and a sharing boundary that protects only the nominal payload. The
format's self-describing structure establishes the mechanism; the Privacy
Exposure Model identifies metadata and structural exposure as privacy threat
classes~\cite{hdf5-data-model,hdf5-privacy-exposure-model}.

\noindent\textbf{Rating basis.}
I2 reflects bounded sensitive-context disclosure; I3 is plausible where the
same mechanism reveals high-sensitivity context or substantially affects a
defined population. E2 reflects the need for both sensitive metadata and a
mis-scoped disclosure boundary. Data sensitivity, affected population, and
workflow prevalence were not measured, so the rating is
\textbf{I2--I3/E2, Moderate--High, Provisional, Moderate confidence}.
Deployment-specific findings should replace this range when those facts are
known. Current installed-population residual risk was not evaluated.

\noindent\textbf{Linked recommendations.}
\par\noindent
\code{CORE-PRIV-S1} and \code{CORE-PRIV-S3}--\code{CORE-PRIV-S4}.

\subsubsection*{CORE-SAFE-001 --- Storage-exhaustion write failure}

\noindent\textbf{Scenario.}
An application writes an HDF5 file to storage that fills during the write,
flush, or close sequence. The operation reports failure through the ordinary
HDF5 error path, including \code{ENOSPC} in the cited example. The library
provides no structured VFD callback at the failure point, no defined recovery
or safe-abandonment procedure, and no supported way to mark the resulting file
as suspect. The file may be incomplete or internally inconsistent, and data
that has not been independently preserved may be lost.

\noindent\textbf{Scope and evidence.}
The assessed population is deployments that write HDF5 files to finite storage
without an assured reserve and that lack a separately verified recoverable copy
of the in-flight data. The controlling evidence is architectural rather than
statistical. In the cited discussion the application receives write and close
errors, while a maintainer states that HDF5 has no recovery story for
\code{ENOSPC} and no callback for the VFD event itself. The discussion records
a workaround --- configuring the metadata cache to discard on evict --- and
enumerates four candidate treatments: space reservation via
\code{posix\_fallocate}, an explicit poisoned-file state, a pre-flush capacity
check, and a virtual-file-driver error callback that lets the caller choose to
retry, fail fast, or panic~\cite{forum-enospc-11714}. General file-corruption
reports provide contextual recurrence evidence but do not establish that
storage exhaustion caused each outcome
~\cite{forum-recover-1146,forum-recover-5441}.

\noindent\textbf{Why this is no longer NE.}
An earlier revision recorded this finding as NE. The additional evidence now
establishes a bounded control deficiency: ordinary operation reports an I/O
failure, but the cited API and maintainer discussion provide no defined
recovery, safe-abandonment, or poisoned-file contract. That is sufficient to
state the scenario and estimate both axes provisionally. It is not evidence
that every \code{ENOSPC} event corrupts a file, that applications receive no
error, or that storage exhaustion is broadly prevalent; those questions remain
part of \code{RM-2D}.

\noindent\textbf{Rating basis.}
I2 reflects a recoverable loss where the workload can be rerun or an
independent copy exists. I3 reflects major loss or corruption where it cannot;
I4 is not assigned because irreversible loss of unique or mission-critical
records is not established by the cited evidence. E2 reflects the material
preconditions: finite storage must fill during an active write, and the
deployment must lack a control that preserves a known-good recoverable state.
The rating is \textbf{I2--I3/E2, Moderate--High, Provisional, Low confidence}.
Confidence is Low because no failure sequence, durable-state result,
recoverability outcome, recurrence rate, or affected population has been
measured; \code{RM-2D} can change the tier as well as confidence. Current
residual risk was not evaluated.

\noindent\textbf{Response.}
Routed Near-term at the upper bound. The Near-term obligation is tested
operational containment --- capacity headroom, free-space monitoring and
alerting, the documented cache workaround where applicable, and an independent
copy of data that cannot be regenerated --- not the library change. The
library treatment in \code{CORE-SAFE-S1} is structural and follows the Planned
and Phase~3 horizons; under Section~\ref{sec:risk-rubric} that split is normal
and does not extend the containment clock.

\noindent\textbf{Linked recommendations.}
\code{CORE-SAFE-S1}; \code{CORE-S12} for the durability metrics.

\noindent\textbf{SSP SIG correspondence.}
\code{HAZ-005}, whose accident chain and safety constraints SC-1 to SC-3
should be adopted as the \code{RM-2D} hypotheses.

\subsubsection*{CORE-SAFE-002 --- Corruption from abrupt termination}

\noindent\textbf{Scenario.}
A writer process is killed, crashes, or loses power while an HDF5 file is open
and metadata updates are in flight. Because in-memory metadata-cache state and
the durable on-disk state can disagree, and because the format has no
journaling or write-ahead log, the file may be left internally inconsistent:
not reopenable, reopenable only after \code{h5clear}, or reopenable but
missing or misrepresenting recently written data. Data written since the last
consistent state may be lost, and the inconsistency may not be detected until a
later read.

\noindent\textbf{Scope and evidence.}
The assessed population is any deployment holding an HDF5 file open for write
across a window in which the process or host can fail --- long acquisitions,
simulation checkpointing, and continuous logging in particular. The evidence is
a recurrence pattern in the community forum spanning sixteen years, from
requests to recover corrupted files and questions about file state after a
crash, through reports that flushing itself sometimes corrupts files, to a 2024
design discussion in which The HDF Group states that a crash or failure ``can
lead to data loss or corruption of the file'' and describes restarting work on
crash-proofing through metadata journaling, write-ahead logging, SWMR-derived
mechanisms, and checkpointing
\cite{forum-recover-1146,forum-crash-state-1598,forum-crash-prevent-2462,
forum-flush-crash-3481,forum-flush-corrupt-4958,forum-recover-5441,
forum-crash-corruption-7581,forum-crashproofing-12003}. That discussion also
records a practitioner report that trapping signals and closing the library
prevents the worst outcomes and usually leaves files intact, with
\code{h5clear} still needed at times --- which is compensating-control
evidence, not a supported guarantee.

\noindent\textbf{Rating basis.}
I2 reflects a recoverable interruption where the file reopens after
\code{h5clear} or the workload can be rerun. I3 reflects major loss or durable
corruption where neither holds. E3 reflects a common workflow --- any
long-running write --- with few material preconditions; abrupt termination
requires no adversary and no unusual configuration. The rating is
\textbf{I2--I3/E3, Moderate--High, Provisional, Low confidence}. Confidence is
Low because the assessment measured no failure injection, established no
version boundaries, and quantified no population: the evidence establishes
recurrence and maintainer acknowledgement, not distribution of outcomes.

\noindent\textbf{Relationship to \code{CORE-SAFE-001}.}
The two are deliberately separate records, because the trigger, the available
containment, and the structural treatment all differ. Storage exhaustion is
foreseeable and can be preflighted; abrupt termination cannot. Rating them
together would obscure both.

\noindent\textbf{Response.}
Routed Near-term at the upper bound, with the same containment-versus-structure
split as \code{CORE-SAFE-001}: the Near-term obligation is flush and
checkpoint discipline, a rehearsed and tested recovery procedure including
\code{h5clear}, and independent preservation of data that cannot be
regenerated. The structural treatment in \code{CORE-SAFE-S2} is a multi-release
programme.

\noindent\textbf{Linked recommendations.}
\code{CORE-SAFE-S2}; \code{CORE-S12} for the durability metrics.

\noindent\textbf{SSP SIG correspondence.}
\code{HAZ-001}, the highest-rated entry in that registry, together with
\code{HAZ-002}, \code{HAZ-006}, \code{HAZ-012}, and \code{HAZ-015}, which
this report does not separately register.

\subsection{Application-Boundary Findings}

\subsubsection*{APP-SEC-001 --- External raw storage confused deputy}

\noindent\textbf{Scenario.}
In the affected July 2026 Hugging Face deployment, a remote party supplied a
valid HDF5 file whose dataset declared local files as external raw storage. The
dataset-processing worker opened those paths with its own filesystem authority
and returned their bytes to the requester, disclosing its environment,
credentials, and source code.

\noindent\textbf{Scope and evidence.}
Hugging Face reports that this vector performed file disclosure without
executing code; the separate Jinja2 injection supplied code execution. The HDF
Group identifies external raw storage as an intentional documented capability
and the service as the authority-bearing confused deputy, not an HDF5 parser
failure~\cite{huggingface-agent-intrusion-2026,hdfgroup-mechanism-not-vulnerability-2026}.
The rating is confined to the affected service at incident time and to
equivalent deployments with the stated preconditions; it is not an
ecosystem-prevalence claim.

\noindent\textbf{Rating basis.}
I3 reflects actual unauthorized resource access and substantial disclosure of
secrets, credentials, and implementation details. E4 reflects normal automatic
processing of an ordinary attacker upload, transparent reference following,
and a direct, repeatable disclosure path within that deployment. The inherent
rating is \textbf{Critical, Rated, High confidence}.

\noindent\textbf{Current residual and response.}
Hugging Face states that its renderer no longer wrongly processes HDF5 external
references. Public evidence is insufficient to validate control design and
coverage, so current residual risk for that service was not evaluated and no
residual tier is assigned. The h5py discussion records laborious application-
level detection, a sandbox recommendation, and the decision to seek upstream
control; HDF5 issue \#6618 records the resulting request for unified controls
over external data, external links, and virtual sources. Proposals discussed
there are not treated as shipped controls
\cite{h5py-gh-issue-2891,hdf5-gh-issue-6618}. Equivalent still-exposed
deployments require Immediate containment.

\noindent\textbf{Linked recommendations.}
\code{APP-S1}, \code{APP-S2}, \code{APP-S4}, and \code{CORE-S4}.

\subsubsection*{APP-SEC-002 --- Keras legacy-HDF5 code deserialization}

\noindent\textbf{Scenario and scope.}
An attacker supplies a crafted legacy \code{.h5}/\code{.hdf5} model to Keras
3.0.0 through 3.11.2. A victim loads it expecting
\code{safe_mode=True} to prevent unsafe reconstruction, but a Lambda layer's
serialized code executes in the loader process
\cite{cve-2025-9905,keras-ghsa-36rr-ww3j-vrjv}.

\noindent\textbf{Rating basis.}
I3 reflects arbitrary code execution in the loader's authority. E2 reflects the
legacy format, crafted Lambda content, affected version, and passive load
preconditions. The rating is \textbf{High, Rated, High confidence}.

\noindent\textbf{Linked recommendations.}
\code{APP-S3}. This is application deserialization layered on HDF5, not a core
HDF5 parser finding.

\subsubsection*{APP-SEC-003 --- Keras external-dataset local-file read}

\noindent\textbf{Scenario and scope.}
A remote attacker supplies a crafted \code{.keras} model to Keras
\code{>=3.0.0,<3.12.1} or \code{>=3.13.0,<3.13.2}. During model loading, an
embedded HDF5 external dataset reference selects a local path and discloses a
file readable by the loader process
\cite{cve-2026-1669,keras-ghsa-3m4q-jmj6-r34q}.

\noindent\textbf{Rating basis.}
I3 reflects unauthorized local-resource access and potentially substantial
sensitive-information disclosure. E3 reflects a remotely supplied artifact,
few technical preconditions, and a normal model-load action, while remaining
bounded to the affected versions and model-loading surface. The rating is
\textbf{High, Rated, High confidence}.

\noindent\textbf{Linked recommendations.}
\code{APP-S1}--\code{APP-S4} and \code{CORE-S4}.

\subsubsection*{APP-SEC-004 --- Keras shape-driven memory exhaustion}

\noindent\textbf{Scenario and scope.}
A remote attacker supplies a crafted \code{.keras} archive to Keras
\code{>=3.0.0,<=3.12.0} or \code{>=3.13.0,<3.13.2}. Its valid
\code{model.weights.h5} declares an extremely large dataset shape; weight
loading allocates without an adequate budget, exhausts memory, and crashes the
Python interpreter~\cite{cve-2026-0897,keras-ghsa-mgx6-5cf9-rr43}.

\noindent\textbf{Rating basis.}
I2 reflects the evidenced process and workflow interruption; no broader outage
or durable data loss is assumed. E3 reflects low-complexity, remotely supplied
content crossing a normal model-loading boundary, subject to the victim or
service loading the artifact. The rating is \textbf{Moderate, Rated, High
confidence}.

\noindent\textbf{Linked recommendations.}
\code{APP-S1}--\code{APP-S3}, \code{CORE-S3}, and \code{CORE-S4}.

\subsubsection*{APP-SEC-005 --- TensorFlow plugin search-path escalation}

\noindent\textbf{Scenario and scope.}
A low-privileged local attacker on a TensorFlow 2.17.0 installation places a
plugin in an unsecured location searched by the HDF5 integration. TensorFlow
loads the plugin and executes its code with the target user's authority
\cite{cve-2026-2492,zdi-26-116,tensorflow-commit-46e7f7f}.

\noindent\textbf{Rating basis.}
I3 reflects privilege escalation and arbitrary code execution. E2 reflects the
required local foothold, affected version, writable searched location, and
plugin-load path. The condition is supported by the coordinated advisory and
vendor fix, but the report did not reproduce it. The rating is
\textbf{High, Rated, Moderate confidence}.

\noindent\textbf{Linked recommendations.}
\code{APP-S5} and \code{CORE-S10}.

\subsection{Independent-Implementation and Interoperability Findings}

\subsubsection*{IND-SAF-001 --- Sufficient nonminimal chunk encoding rejected}

\noindent\textbf{Scenario.}
A workflow upgrades its reader to an HDF5 2.0/2.1 release containing the
exact-width check---confirmed in 2.0.0, 2.1.0, and 2.1.1---and opens a
version-4 chunk layout whose encoded dimension width is sufficient but
nonminimal. Validation rejects the file even though the format did not require
the smallest possible width, interrupting access to data that earlier readers
accepted.

\noindent\textbf{Scope and evidence.}
HDF5 issue \#6409 distinguishes the specification-permitted encoding from an
undersized invalid encoding and records effects on independent writers. The
merged reader-side change accepts any sufficient width and was tracked against
the HDF5 2.2.0 milestone; this report does not claim that a specific release
artifact or backport contains it without release-level verification
\cite{hdf5-gh-issue-6409,hdf5-gh-pr-6432,purehdf-gh-issue-168,purehdf-pr-160}.

\noindent\textbf{Rating basis.}
I2 reflects a meaningful, bounded interoperability and data-access failure. E2
reflects the specific layout version, nonminimal encoding, and reader-upgrade
preconditions. The rating is \textbf{Moderate, Rated, High confidence} for the
historical affected scope. Current residual risk was not evaluated, and no
residual tier is assigned.

\noindent\textbf{Linked recommendations.}
\code{CORE-S1}, \code{CORE-S6}, \code{CORE-S8}, and \code{IND-S1}.

\subsubsection*{IND-SAF-002 --- Contradictory PureHDF chunk metadata}

\noindent\textbf{Scenario.}
A production or archival workflow upgrades from a tolerant HDF5 1.x reader to
HDF5 2.0 and reads an affected compressed file produced by PureHDF 2.1.2. The
file's version-4 chunk layout contradicts its datatype metadata; the stricter
reader correctly rejects the affected datasets, interrupting processing and
stranding existing data until the reader is pinned or the files are repaired or
regenerated.

\noindent\textbf{Scope and evidence.}
In the supplied production-derived file, the HDF5 maintainer analysis found
float64 datatype element size 8 while the chunk layout recorded stored element
size 4. Older HDF5 1.10.10 and 1.14.6 readers accepted the bytes; HDF5 2.0
rejected them. The issue was closed as malformed writer output, and the PureHDF
discussion records the writer-side response~\cite{hdf5-gh-issue-6534,purehdf-gh-issue-164}.
This condition is distinct from the sufficient nonminimal width in
\code{IND-SAF-001}; permissive acceptance is not evidence of conformance.

\noindent\textbf{Rating basis.}
I2 reflects actual production workflow interruption and a persistent but
repairable compatibility problem; no irreversible data loss is evidenced. E2
reflects the confirmed writer version, affected layout, compressed dataset, and
reader-upgrade preconditions. The rating is \textbf{Moderate, Rated, High
confidence}. A technical metadata-repair path was proposed and older readers
remained a workaround, although the reporter said rewriting affected production
files was not operationally feasible in that pipeline. Later PureHDF output
restored production operation, but the inventory and repair status of existing
files are unknown. Current residual risk for the legacy corpus was not
evaluated, and no residual tier is assigned.

\noindent\textbf{Linked recommendations.}
\code{CORE-S1}, \code{CORE-S6}, \code{CORE-S8}, and \code{IND-S1}.

\subsection{Format and Specification Evidence Record}

\subsubsection*{FMT-SAF-001 --- Unresolved normative interpretation}

\noindent\textbf{Observation and scope.}
Issue~\#6616 records a specific unresolved disagreement. The published format
describes the group-info message as optional, while the C library reports
``message type not found'' when it is absent and therefore cannot modify the
file. The reporter identifies two possible resolutions: the library supplies
the documented defaults, or the specification is corrected to require the
message. At the evidence cutoff the issue was open, assigned, labelled as a
documentation bug, and milestoned, with no authoritative ruling recorded
~\cite{hdf5-gh-issue-6616}.

Issue~\#6409 is deliberately not used as a second instance of indeterminacy.
There the published width rule was sufficient to decide conformance, the
reference reader was wrong, and the reader was corrected
~\cite{hdf5-gh-issue-6409,hdf5-gh-pr-6432}. It remains strong evidence that an
independent implementation can expose a reference-library defect; it does not
show that the normative artifact was unable to settle the question.

\noindent\textbf{Why this remains in the evidence backlog.}
The issue establishes a governance and specification-adjudication gap, but the
available record does not identify a bounded file population, a material user
or data consequence, or recurrence of the same unresolved condition. The
separate \code{IND-SAF-001} and \code{IND-SAF-002} findings already rate the
evidenced consequences of their particular reader/writer divergences. Combining
those outcomes with issue~\#6616 would rate heterogeneous mechanisms twice and
would use a resolved case to support an unresolved one. Impact, exposure, tier,
and risk-response urgency are therefore \textbf{NE}; NE is not a Low-risk
conclusion.

\noindent\textbf{Evidence and governance action.}
Within 90 days of roadmap adoption, the format authority should adjudicate
issue~\#6616, publish the rationale and authoritative artifact, change whichever
artifact is wrong, and preserve positive and negative fixtures. That is an
administrative assurance target, not an inherent-risk response clock. To
promote this record, identify files or workflows that depend on the disputed
rule and measure the resulting access, integrity, or compatibility consequence.
The systematic prevention work remains \code{FMT-S1}: a standing adjudication
procedure and a machine-readable description for a declared subset, with
implementation-defined behavior permitted only when it is explicit and tested.

\noindent\textbf{Linked recommendations.}
\code{FMT-S1}; \code{IND-S1} for paired fixtures; and \code{CORE-S1}, whose
invariant registry must distinguish specification-derived requirements,
implementation-defined behavior, and reference-implementation behavior.

\noindent\textbf{SSP SIG correspondence.}
No direct counterpart. The hazard model has a parser-ambiguity family
(\code{H6}) and the independent-implementation model has \textbf{I1}
parser/spec divergence, but neither registry contains an entry for an unresolved
normative-artifact question.

\subsection{Long-Term Accessibility Findings}

\subsubsection*{LONG-SAF-001 --- Filter-dependent loss of interpretability}

\noindent\textbf{Scenario.}
An HDF5 file is retained in an archive and a dataset in it requires a non-core
filter plugin. At some later point the filter is absent from the reader's build
or platform, unbuildable against the current toolchain, or no longer
obtainable. The file bytes survive, but the affected dataset cannot be decoded.

\noindent\textbf{Assessment horizon.}
This finding states an assessment horizon of \textbf{ten years or more} from
the date of writing. The cited reports establish that missing-filter failures
already occur; they do not measure how extension availability changes over a
decade. The horizon is therefore a material assumption and, under decision
rule 7 of Section~\ref{sec:risk-rubric}, this rating must not be directly
ordered against ratings that state no horizon.

\noindent\textbf{Scope and evidence.}
The assessed population is archived files containing datasets that require a
non-core filter. The mechanism is observable at short horizons: community
reports record users unable to read data because a required compression filter
was absent from their build or platform, across a fifteen-year span and several
filter families
\cite{forum-filter-compressed-2756,forum-filter-h5dump-3396,
forum-filter-dynamic-4047,forum-filter-lzo-4102}. The archival extrapolation
--- that availability decreases rather than increases over a decade --- is
reasoned inference, not measurement, and is the reason confidence is Low.

VFD and VOL dependencies are excluded from the rated population. A missing VFD
can affect whether bytes can be reached at all, while a VOL connector may map
to a different store rather than to a self-contained HDF5 file. Those are
materially different mechanisms and affected populations, and the report has
no comparable archival failure evidence for them. They remain reasoned
extensions of the analysis and require separate scenarios before rating.

\noindent\textbf{Rating basis.}
I2 reflects the evidenced material workflow interruption: an affected dataset
cannot be read until the filter is recovered, rebuilt, or re-implemented. I3
would require evidence that unavailable decoding caused consequential or
irrecoverable loss of significant records; the cited reports do not establish
that outcome. E2 reflects the material preconditions --- an archived dataset
must depend on a non-core filter and that filter must be unavailable to the
reader. The rating is \textbf{I2/E2, Moderate, Provisional, Low confidence}, at
the stated horizon. Confidence is Low because the archived population,
dependency survival distribution, and decadal availability trend were not
measured.

\noindent\textbf{Response.}
The Moderate rating routes \emph{Planned}. Within 90 days, the accountable
archive or producer should define a bounded corpus, inventory each file's
filter dependencies, and record whether each dependency is available,
reproducible, removable by migration, or accepted with an owner and review
date. The longer-term archival profile and preserved specifications are
producer-side and curator-side work delivered through \code{LONG-S1}; they do
not defer that bounded first deliverable.

\noindent\textbf{Linked recommendations.}
\code{LONG-S1}; \code{IND-S1} for the fixture and reader-matrix machinery it
reuses; \code{CORE-S11} and \code{APP-S6} for the provenance inputs that make a
dependency discoverable.

\noindent\textbf{SSP SIG correspondence.}
\code{HAZ-014} covers the same filter-plugin portability mechanism at Moderate.
Neither registry covers VFD or VOL archival dependency; this report identifies
those as unmeasured extensions rather than including them in this rating.

\subsection{Evidence-Backlog Record Details}

\subsubsection*{EXT-PRIV-001 --- Extension-mediated privacy boundary observation}

\noindent\textbf{Observation and evidence gap.}
VOL connectors, filters, and VFDs can move data or metadata into remote stores,
caches, temporary artifacts, logs, networks, or dynamically loaded code and
thereby change the privacy boundary described in
Section~\ref{sec:extensions}. The report does not define a concrete deployment,
sensitive-data population, unauthorized recipient, retention period, or harm
magnitude for this general scenario~\cite{hdf5-privacy-exposure-model}.

\noindent\textbf{Register treatment.}
This mechanism is not a formal finding. Impact, exposure, tier, confidence, and
response urgency are not assigned; this does not dispute the documented
mechanism.

\noindent\textbf{Linked recommendations.}
\code{CORE-PRIV-S2}, \code{CORE-PRIV-S3}, \code{VOL-S1}--\code{VOL-S4},
\code{FILTER-S1}--\code{FILTER-S5}, and \code{VFD-S1}--\code{VFD-S4}.

\subsubsection*{SUP-SAFE-001 --- Downstream build and feature mismatch observation}

\noindent\textbf{Observation and evidence gap.}
A downstream environment loads an HDF5 build whose ABI, linkage, filters,
plugins, VFDs, parallel features, or file-locking behavior differs from the
profile assumed by a file or workflow, causing a read failure or operational
interruption. Section~\ref{sec:supply-chain} establishes the breadth and
configuration mechanisms but does not establish an affected population,
frequency, bounded version set, or realized impact.

\noindent\textbf{Register treatment.}
This theme is not a formal finding. Impact, exposure, tier, confidence, and
response urgency are not assigned; a material workflow interruption remains a
plausible hypothesis, not a report-wide rating.

\noindent\textbf{Linked recommendations.}
\code{CORE-S11}, \code{APP-S6}, and \code{SUP-S1}.

\end{document}
